\section{指揮デバイスのプログラムソース}
\label{sec:apndx_siki}
本節では，指揮デバイスの制御プログラムの全文を示す．ファイル構成は\ref{sec:siki_naibu}節の表\ref{table:siki_files}の通りである．

\subsection{Controller\_simultaneous.ino}
\begin{lstlisting}[caption={Controller\_simultaneous.ino}, label={lst:siki_ino}]
#include "SyncManager.h"
#include "PotManager.h"
#include "AccManager.h"

#include <Wire.h>
#include <math.h>

// Wi-Fi のネットワーク名(hackathon001-WPA2)
const char SSID[] = "hackathon001-WPA2";
// その Wi-Fi のパスワード(hackathon001)
const char PASS[] = "hackathon001";

const uint8_t input_PIN = 3; // スイッチの状態を読み取るピン
bool switch_status;// ボタンの状態を保持
bool before_switch_status = 0;// 前回のボタンの状態を保持
bool isButtonPlaying = false;// 開始・停止信号を交互に切り替えるフラグ

SyncManager sync;
PotManager pot;
AccManager acc;

void setup() {
    Serial.begin(115200);
    while (!Serial) {}

    Serial.println("--- Setup Started ---");
    pinMode(input_PIN, INPUT);

    Serial.println("Attempting to connect to WiFi...");
    sync.begin(SSID, PASS);

    Serial.println("\nWiFi Connected successfully!");
    Serial.println("Commands: button -> Start/End");

    acc.setupADXL345();

    before_switch_status = digitalRead(input_PIN);
}

void loop() {
    button();
    pot.update();
    acc.update(sync, pot);
}

// 物理ボタンの押下を検知し，開始信号・停止信号のトグル送信と演奏状態フラグの反転を一括管理する．
void button() {
    switch_status = digitalRead(input_PIN);
    if (switch_status == 0 && before_switch_status == 1) {
        if (isButtonPlaying) {
            Serial.println("Button: [END] sending UDP...");
            sync.sendEnd();
            isButtonPlaying = false;
        } else {
            Serial.println("Button: [START] sending UDP to all simultaneously...");
            uint8_t step = pot.getBpmStep();
            sync.sendStart(step);
            isButtonPlaying = true;
        }
        delay(300); // チャタリング防止
    }
    before_switch_status = switch_status;
}
\end{lstlisting}

\subsection{AccManager.h}
\begin{lstlisting}[caption={AccManager.h}, label={lst:siki_acc_h}]
// 加速度センサー管理クラス
#ifndef ACC_MANAGER_H
#define ACC_MANAGER_H

#include <Arduino.h>
#include <Wire.h>
#include "SyncManager.h"
#include "PotManager.h"

class AccManager {
public:
    AccManager();
    // ADXL345（加速度センサー）のI2C初期化および計測開始の準備を行う．
    void setupADXL345();  
    // 20ms周期の時間管理と，閾値判定による指揮棒のシェイク動作検知を行う．
    void update(SyncManager& sync, PotManager& pot); 
    // potの予約フラグを確認し，変化があればsyncを通じてBPMや音量の段階番号を送信する．
    void updateShakedInfo(SyncManager& sync, PotManager& pot);

private:
    float readAccMagnitude(); // I2Cで読み取ったXYZの3軸生データから合成ベクトル(加速度)を算出
    void  recoverI2C(); // I2Cバスロックアップ(切断やノイズなどによりI2C通信が不能になったとき)からの復旧を行う．

    const uint8_t ADXL345_ADDR = 0x53; // 加速度センサー(ADXL345)のI2Cアドレスを定義
    const uint8_t REG_POWER_CTL = 0x2D; // 加速度センサの電源管理をするスイッチを定義
    const uint8_t REG_DATA_FORMAT = 0x31; // 測定範囲設定用レジスタのアドレスを定義
    const uint8_t REG_DATAX0 = 0x32; // 測定した値が書き込まれる場所を定義

    const float accThreshold = 110.0f; // 振ったと判断するための変化量の閾値を定義
    const unsigned long shakeInterval = 750; // 連続検知を防止するための待機時間を定義
    const unsigned long ACC_SAMPLE_INTERVAL = 20;// 加速度センサのサンプリング周期を定義

    // 内部変数
    float currentAccVal = 0.0f; // 加速度センサからの値を格納
    float lastAccVal = 0.0f;  // 変化量算出のための直前の加速度センサ値を格納
    unsigned long lastDetectTime = 0;  // 前回振ったことを検知した時間をms単位で格納
    unsigned long lastSampleTime = 0;// 加速度センサの読み取り間隔を管理するための前回の時刻を格納
    unsigned long lastRecoverMs  = 0;    // 復旧ログ処理の連続実行を防ぐための時刻を保持

    // 振られたときに反転
    bool isStarted = false; 
};

#endif
\end{lstlisting}

\subsection{AccManager.cpp}
\begin{lstlisting}[caption={AccManager.cpp}, label={lst:siki_acc_cpp}]
#include "AccManager.h"

AccManager::AccManager(){
  // コンストラクタ（初期化処理なし）
}

// ADXL345（加速度センサー）のI2C初期化および計測開始の準備を行う．
void AccManager::setupADXL345() {
  Wire.begin();

  Wire.beginTransmission(ADXL345_ADDR);
  Wire.write(REG_POWER_CTL);
  Wire.write(0x08); // Measure bit ON
  Wire.endTransmission();

  Wire.beginTransmission(ADXL345_ADDR);
  Wire.write(REG_DATA_FORMAT);
  Wire.write(0x00);
  Wire.endTransmission();

  lastAccVal = readAccMagnitude();
  
  Serial.println("GY-291 初期化完了");
  Serial.print("閾値: "); 
  Serial.print(accThreshold);
  Serial.print("インターバル: "); 
  Serial.print(shakeInterval); 
  Serial.println(" ms");
}

// 20ms周期の時間管理と，閾値判定による指揮棒のシェイク動作検知を行う．
void AccManager::update(SyncManager& sync, PotManager& pot) {
  unsigned long now = millis();
  
  if (now - lastSampleTime >= ACC_SAMPLE_INTERVAL) {
    lastSampleTime = now;

    currentAccVal = readAccMagnitude();
    float diff = fabsf(currentAccVal - lastAccVal);

    bool overThreshold  = (diff > accThreshold);
    bool intervalPassed = ((now - lastDetectTime) >= shakeInterval);

    // 500msごとにdiff値を表示（閾値チューニング用）
    static unsigned long lastDebugTime = 0;
    if (now - lastDebugTime >= 500) {
      lastDebugTime = now;
      //Serial.print("[ACC] diff="); Serial.print(diff, 1);
      //Serial.print(" / 閾値="); Serial.println(accThreshold, 1);
    }

    if (overThreshold && intervalPassed) {
      lastDetectTime = now;
      isStarted = !isStarted;

      Serial.print("[SHAKE検知] diff=");
      Serial.print(diff, 4);
      Serial.println(isStarted ? " | START送信" : " | END送信");

      updateShakedInfo(sync, pot);
    }

    lastAccVal = currentAccVal;
  }
}

// I2Cで読み取ったXYZの3軸生データから合成ベクトル(加速度)を算出する．
float AccManager::readAccMagnitude() {
  Wire.beginTransmission(ADXL345_ADDR);
  Wire.write(REG_DATAX0);
  // endTransmission!=0 はNACK/バス異常。リピーテッドスタート前に検出する。
  if (Wire.endTransmission(false) != 0) {
    recoverI2C();
    return lastAccVal;   // 直前値を返してニセshakeを防ぐ
  }

  uint8_t n = Wire.requestFrom(ADXL345_ADDR, (uint8_t)6, (uint8_t)true);
  if (n < 6) {             // 6バイト読めない＝ロックアップ等
    recoverI2C();
    return lastAccVal;
  }

  int16_t rawX = (int16_t)((Wire.read()) | (Wire.read() << 8));
  int16_t rawY = (int16_t)((Wire.read()) | (Wire.read() << 8));
  int16_t rawZ = (int16_t)((Wire.read()) | (Wire.read() << 8));

  float fx = (float)rawX;
  float fy = (float)rawY;
  float fz = (float)rawZ;

  return sqrtf(fx * fx + fy * fy + fz * fz);
}

// I2Cバスロックアップ(切断やノイズなどによりI2C通信が不能になったとき)からの復旧を行う．
void AccManager::recoverI2C() {
  unsigned long now = millis();
  if (now - lastRecoverMs >= 1000) {   // ログのスロットル（毎20msの連発を防ぐ）
    lastRecoverMs = now;
    Serial.println("[ACC] I2C異常検出 → バス復旧を実行");
  }

  Wire.end();

  // バスクリア: SCLを9クロック叩いて握られたSDAを解放
  pinMode(SDA, INPUT_PULLUP);
  pinMode(SCL, OUTPUT);
  for (int i = 0; i < 9; i++) {
    digitalWrite(SCL, LOW);  delayMicroseconds(5);
    digitalWrite(SCL, HIGH); delayMicroseconds(5);
  }
  // STOP条件 (SDA: LOW→HIGH while SCL HIGH)
  pinMode(SDA, OUTPUT);
  digitalWrite(SDA, LOW);  delayMicroseconds(5);
  digitalWrite(SCL, HIGH); delayMicroseconds(5);
  digitalWrite(SDA, HIGH); delayMicroseconds(5);

  // Wire再初期化＆センサ再設定（setupADXL345の通信部と同じ）
  Wire.begin();
  Wire.beginTransmission(ADXL345_ADDR);
  Wire.write(REG_POWER_CTL);   Wire.write(0x08);
  Wire.endTransmission();
  Wire.beginTransmission(ADXL345_ADDR);
  Wire.write(REG_DATA_FORMAT); Wire.write(0x00);
  Wire.endTransmission();
}

// potの予約フラグを確認し，変化があればsyncを通じてBPMや音量の段階番号を送信する．
void AccManager::updateShakedInfo(SyncManager& sync, PotManager& pot) {
  if (pot.bpmFrag) {
    sync.sendBPM(pot.getBpmStep());
    Serial.print("[SHAKE] BPM送信 step="); Serial.println(pot.getBpmStep());
    pot.bpmFrag = false;
  }
  if (pot.volFrag) {
    sync.sendVolume(pot.getVolStep());
    Serial.print("[SHAKE] VOL送信 step="); Serial.println(pot.getVolStep());
    pot.volFrag = false;
  }
}
\end{lstlisting}

\subsection{PotManager.h}
\begin{lstlisting}[caption={PotManager.h}, label={lst:siki_pot_h}]
// 可変抵抗器管理クラス
#ifndef POT_MANAGER_H
#define POT_MANAGER_H

#include <Arduino.h>
#include "SyncManager.h"

class PotManager {
public:
    PotManager();
    
    // 20ms周期の時間管理を行い，可変抵抗器の値のサンプリング処理を駆動する．
    void update();
    // カプセル化された現在のBPM段階番号（1〜5）を外部に返す．
    uint8_t getBpmStep() const { return currentBpmStep; }
    // カプセル化された現在の音量段階番号（1〜5）を外部に返す．
    uint8_t getVolStep() const { return 6 - currentVolStep; }

    
    bool bpmFrag = false;// 値の変化時に指揮棒の振りと連動させるためのBPM送信予約フラグ
    bool volFrag = false;// 値の変化時に指揮棒の振りと連動させるための音量送信予約フラグ

private:
    // 平滑化された平均アナログ値を，設計書に基づく境界値に照らし合わせて1〜5の段階番号に変換する．
    int calcStep(float val);
    // BPM用可変抵抗器の値を移動平均で処理し，前回値からの変化を検知してフラグ（真偽値）を返す．
    bool bpmUpdatePotentiometers();
    // 音量用可変抵抗器の値を移動平均で処理し，前回値からの変化を検知してフラグ（真偽値）を返す．
    bool volUpdatePotentiometers();

    // ピン定義
    const int PIN_BPM = A0;  // BPM変更用の可変抵抗器を接続するアナログ入力ピン番号を定義
    const int PIN_VOL = A1; // 音量変更用の可変抵抗器を接続するアナログ入力ピン番号を定義
    const unsigned long SAMPLE_INTERVAL = 20; // 可変抵抗器の値を読み取るサンプリング周期
    const int sampleSize = 10; // 移動平均算出のためのサンプル数を定義
    const int STEP_BOUNDARIES[4] = {300, 500, 650, 818};// アナログ入力値を5段階に変換するための閾値

    // 内部変数
    unsigned long lastSampleTime = 0; // 前回のサンプリング時刻
    
    int   bpmSamples[10] = {0}; // BPM用のサンプル値を保持するための配列
    long  bpmTotal = 0; // bpmSamples内の累積合計を格納
    float averageBpmVal = 0.0f; // 算出されたBPM用のアナログ入力の平均値を格納
   
    int   volSamples[10] = {0};  // 音量用のサンプル値を保持するための配列
    long  volTotal = 0; // volSamples内の累積合計を格納
    float averageVolVal = 0.0f; // 算出された音量用のアナログ入力の平均値を格納
    
    int  bpmReadIndex = 0;// BPMの配列内での現在の書き込み位置を管理するインデックス
    int  volReadIndex = 0;// Volumeの配列内での現在の書き込み位置を管理するインデックス
    bool bpmBufferFull = false;// BPM用の移動平均バッファが初回サンプルで満たされたかを判定するフラグ
    bool volBufferFull = false;// 音量用の移動平均バッファが初回サンプルで満たされたかを判定するフラグ

    // 算出されたBPMの5段階のレベル（1～5）を格納
    uint8_t currentBpmStep = 0; 
    // 算出された音量の5段階のレベル（1～5）を格納
    uint8_t currentVolStep = 0; 
    
    int prevBpmStep = 0;// 前回のBPM段階を保持
    int prevVolStep = 0;// 前回の音量段階を保持
};

#endif
\end{lstlisting}

\subsection{PotManager.cpp}
\begin{lstlisting}[caption={PotManager.cpp}, label={lst:siki_pot_cpp}]
#include "PotManager.h"

PotManager::PotManager(){
  // コンストラクタ（初期化処理なし）
}

// 20ms周期の時間管理を行い，可変抵抗器の値のサンプリング処理を駆動する．
void PotManager::update() {
  unsigned long now = micros();
  if (now - lastSampleTime >= SAMPLE_INTERVAL) {
    lastSampleTime = now;
    if (bpmUpdatePotentiometers()) bpmFrag = true;
    if (volUpdatePotentiometers()) volFrag = true;
  }
}

// 平滑化された平均アナログ値を，設計書に基づく境界値に照らし合わせて1〜5の段階番号に変換する．
int PotManager::calcStep(float val) {
  if      (val <= STEP_BOUNDARIES[0]) return 1;
  else if (val <= STEP_BOUNDARIES[1]) return 2;
  else if (val <= STEP_BOUNDARIES[2]) return 3;
  else if (val <= STEP_BOUNDARIES[3]) return 4;
  else                                return 5;
}

// BPM用可変抵抗器の値を移動平均で処理し，前回値からの変化を検知してフラグ（真偽値）を返す．
bool PotManager::bpmUpdatePotentiometers() {
  bpmTotal -= bpmSamples[bpmReadIndex];
  bpmSamples[bpmReadIndex] = analogRead(PIN_BPM);
  bpmTotal += bpmSamples[bpmReadIndex];

  bpmReadIndex = (bpmReadIndex + 1) % sampleSize;
  if (bpmReadIndex == 0) bpmBufferFull = true;
  if (!bpmBufferFull) return false;

  averageBpmVal = (float)bpmTotal / sampleSize;
  currentBpmStep = calcStep(averageBpmVal);

  if (currentBpmStep != prevBpmStep) {
    Serial.println("------ BPMパラメータ更新 ------");
    Serial.print("[BPM] 平均センサ値="); Serial.print(averageBpmVal, 1);
    Serial.print("  段階=");          Serial.print(currentBpmStep);
    prevBpmStep = currentBpmStep;
    return true;
  }
  return false;
}

// 音量用可変抵抗器の値を移動平均で処理し，前回値からの変化を検知してフラグ（真偽値）を返す．
bool PotManager::volUpdatePotentiometers() {
  volTotal -= volSamples[volReadIndex];
  volSamples[volReadIndex] = analogRead(PIN_VOL);
  volTotal += volSamples[volReadIndex];

  volReadIndex = (volReadIndex + 1) % sampleSize;
  if (volReadIndex == 0) volBufferFull = true;
  if (!volBufferFull) return false;

  averageVolVal = (float)volTotal / sampleSize;
  currentVolStep = calcStep(averageVolVal);

  if(currentVolStep != prevVolStep){
    Serial.println("------ VOLパラメータ更新 ------ ");
    Serial.print("[VOL] 平均センサ値="); Serial.print(averageVolVal, 1);
    Serial.print("  段階=");          Serial.print(6 - currentVolStep);
    prevVolStep = currentVolStep;
    return true;
  }
  return false;
}
\end{lstlisting}

\subsection{SyncManager.h}
\begin{lstlisting}[caption={SyncManager.h}, label={lst:siki_sync_h}]
#ifndef SYNC_MANAGER_H
#define SYNC_MANAGER_H

#include <WiFiS3.h>
#include <WiFiUdp.h>

extern IPAddress fluteIP;// フルートのIPアドレス
extern IPAddress pianoIP;// ピアノのIPアドレス
extern IPAddress stringsIP;// ストリングスのIPアドレス
extern IPAddress drumIP;// ドラムのIPアドレス
extern IPAddress ferrisWheelIP;// 観覧車（送信）のIPアドレス
extern IPAddress ferrisWheelGroup; // 観覧車（送信）のIPアドレス(マルチキャスト用)
extern IPAddress instrumentGroup;  // 楽器全体のIPアドレス

class SyncManager {
public:
    // コンストラクタ（ポート番号などを初期化する）
    SyncManager();
    
    // setup内で一度だけ呼び出される関数であり，Wi-FiとUDPを準備している．
    void begin(const char ssid[], const char pass[]);
    
    // 演奏開始合図（'S'）とBPMの段階番号をラウンドロビン(順次割り当て)で送信．
    void sendStart(uint8_t stepNumber);
    
    // BPM情報を観覧車へユニキャスト送信するための関数．
    void sendBPM(uint8_t stepNumber);
    
    // 音量情報を各楽器機へラウンドロビンでユニキャスト送信するための関数．
    void sendVolume(uint8_t stepNumber);
    
    // 演奏終了合図（'E'）を全機へラウンドロビンで送信する関数．
    void sendEnd();

private:
    // ラウンドロビン方式で到達時間差を無くし，かつパケットロス対策として3回冗長送信を行うための共通処理．
    // （※ターゲットごとに3連射してから次のターゲットへ進む方式だと、後方の
    // ターゲットほど到達が遅れて機器間に無視できない時間差が生まれるため、
    // 1周ごとに全ターゲットへ送ってから次の周に進む方式にしている。）
    void packetRoundRobin(const uint8_t packet[], uint8_t b, const IPAddress targets[], uint8_t targetCount);
};

#endif
\end{lstlisting}

\subsection{SyncManager.cpp}
\begin{lstlisting}[caption={SyncManager.cpp}, label={lst:siki_sync_cpp}]
#include "SyncManager.h"


const uint16_t TARGET_PORT = 5005;// ターゲットとなるポート番号
const uint16_t LOCAL_PORT = 8080;// 自身のポート番号
WiFiUDP Udp;


IPAddress fluteIP(192, 168, 11, 14);// フルートのIPアドレス
IPAddress pianoIP(192, 168, 11, 12);// ピアノのIPアドレス
IPAddress stringsIP(192, 168, 11, 13);// ストリングスのIPアドレス
IPAddress drumIP(192, 168, 11, 11);// ドラムのIPアドレス
IPAddress ferrisWheelIP(192, 168, 11, 10);   // 観覧車（送信）のIPアドレス
IPAddress ferrisWheelGroup(239, 255, 0, 1);  // 観覧車（送信）のIPアドレス(マルチキャスト用) ※未使用
IPAddress instrumentGroup(239, 255, 0, 2);// 楽器全体のIPアドレス

SyncManager::SyncManager() {
    // コンストラクタ（今回は初期化処理なし）(インスタンスが作られた瞬間に，自動的に一番最初に1回だけ実行される特殊な関数)
}

// setup内で一度だけ呼び出される関数であり，Wi-FiとUDPを準備している．
void SyncManager::begin(const char ssid[], const char pass[]) {
    WiFi.begin(ssid, pass); // Wi-Fi接続処理を開始
    // 接続状態が「WL_CONNECTED（接続完了）」になるまでループ待機
    while (WiFi.status() != WL_CONNECTED) {
    }
    Udp.begin(LOCAL_PORT); // 自身のポートを開放し、UDP通信を有効化
}

// 演奏開始合図（'S'）とBPMの段階番号をラウンドロビン(順次割り当て)で送信．
void SyncManager::sendStart(uint8_t stepNumber) {
    // 設計書通りの2バイト固定長ペイロード（第1バイト:ヘッダ'S', 第2バイト:BPM段階番号）
    uint8_t startPacket[2] = {'S', stepNumber};
    // ドラムはisInstJoinを'S'受信時にリセットするため、他機より早く
    // 届かせたい。ただしラウンドロビン送信なら全ターゲットへの到達は
    // ほぼ同時になるため、この並び自体はもう重要ではない。
    IPAddress targets[] = {drumIP, ferrisWheelIP, pianoIP, stringsIP, fluteIP};
    packetRoundRobin(startPacket, 2, targets, 5);
}

// BPM情報を観覧車へユニキャスト送信するための関数．
void SyncManager::sendBPM(uint8_t stepNumber){
    uint8_t bpmPacket[2] = {'B',stepNumber};
    IPAddress targets[] = {ferrisWheelIP};
    packetRoundRobin(bpmPacket, 2, targets, 1);
}

// 音量情報を各楽器機へラウンドロビンでユニキャスト送信するための関数．
void SyncManager::sendVolume(uint8_t stepNumber){
    uint8_t volumePacket[2] = {'V',(uint8_t)(6 - stepNumber)};
    IPAddress targets[] = {fluteIP, pianoIP, stringsIP, drumIP};
    packetRoundRobin(volumePacket, 2, targets, 4);
}

// 演奏終了合図（'E'）を全機へラウンドロビンで送信する関数．
void SyncManager::sendEnd(){
    uint8_t endPacket[2] = {'E', 0x00};
    IPAddress targets[] = {drumIP, ferrisWheelIP, pianoIP, stringsIP, fluteIP};
    packetRoundRobin(endPacket, 2, targets, 5);
}

// ラウンドロビン方式で到達時間差を無くし，かつパケットロス対策として3回冗長送信を行うための共通処理．
// 1周目で全ターゲットに1回ずつ送り、20ms空けてから2周目、3周目と繰り返す。
// ターゲットごとに3連射してから次へ進む方式だと、後方のターゲットほど
// 到達が遅れて機器間で無視できない時間差が生まれるため、この方式にしている。
void SyncManager::packetRoundRobin(const uint8_t packet[], uint8_t b, const IPAddress targets[], uint8_t targetCount){
    for (int round = 0; round < 3; round++) {
        for (uint8_t t = 0; t < targetCount; t++) {
            Udp.beginPacket(targets[t], TARGET_PORT); // 送信先IPとポートをセット
            Udp.write(packet, b);
            Udp.endPacket(); // パケットをパースして実際に送信
        }
        // 最後の周（3周目）以外は，20msのインターバル（遅延）を挟む
        if (round < 2) {
            delay(20);
        }
    }
}
\end{lstlisting}
